\documentclass[../../../include/open-logic-section]{subfiles}

\begin{document}

\olfileid[id]{sth}{replacement}{absinf}
\olsection{Penggantian dan ``Ketakhinggaan Mutlak''}

Sekarang kita akan meninggalkan \limofsize{} dan menjelajahi keluarga upaya
(intrinsik) lain untuk membenarkan Penggantian, yang benar-benar menganggap
serius gagasan bahwa himpunan dibentuk dalam tahap-tahap.

Ketika pertama kali menguraikan proses iteratif, kita mengajukan beberapa
prinsip yang menjelaskan apa yang terjadi pada setiap tahap. Prinsip-prinsip
ini ialah \stageshier, \stagesord, dan \stagesacc. Kemudian kita menambahkan
beberapa prinsip yang memberi tahu kita sesuatu mengenai banyaknya tahap:
\stagessucc{} memberi tahu kita bahwa proses pembentukan himpunan tidak
pernah berakhir, dan \stagesinf{} memberi tahu kita bahwa proses itu melewati
suatu tahap yang berindeks tak hingga.

Masuk akal untuk mengusulkan bahwa kedua prinsip terakhir ini mengikuti
suatu prinsip yang lebih luas, seperti:
\begin{enumerate}
	\item[] \stagesinex. Banyaknya tahap bersifat tak hingga secara mutlak;
	hierarki itu setinggi mungkin.
\end{enumerate}
Jelas bahwa ini merupakan prinsip informal. Namun, sekalipun prinsip itu
tidak langsung \emph{merupakan konsekuensi} dari konsepsi kumulatif-iteratif tentang
himpunan, prinsip tersebut tentu tampak \emph{selaras} dengannya. Setidaknya,
dan tidak seperti \limofsize, prinsip ini mempertahankan gagasan bahwa
himpunan dibentuk tahap demi tahap.

Sekarang harapannya ialah memanfaatkan \stagesinex{} menjadi pembenaran bagi
Penggantian. Jadi, mari kita lihat bagaimana hal ini dapat dilakukan.

Dalam \olref[ordinals][idea]{sec}, kita melihat bahwa mudah untuk
mengonstruksi suatu pengurutan baik yang (secara moral) seharusnya isomorfik
dengan $\omega+\omega$. Dengan kata lain, kita dapat dengan mudah
membayangkan proses iteratif tahap demi tahap yang tipe urutannya (secara
moral) ialah $\omega+\omega$. Karena itu, jika kita telah menerima
\stagesinex, tentu kita juga seharusnya menerima bahwa setidaknya terdapat
tahap ke-$\omega+\omega$ dalam hierarki, yakni $V_{\omega+\omega}$, sebab
hierarki itu tentu \emph{dapat} berlanjut sejauh ini.

Gagasan ini dapat diperumum sebagai berikut: untuk setiap pengurutan baik,
proses pembangunan hierarki iteratif seharusnya berjalan setidaknya sejauh
pengurutan baik itu. Dan kita dapat menjaminnya hanya dengan memperlakukan
\olref[ordinals][ordtype]{thmOrdinalRepresentation} sebagai suatu
\emph{aksioma}. Hal ini akan memberi tahu kita bahwa setiap pengurutan baik
isomorfik dengan suatu ordinal von Neumann. Karena setiap ordinal von Neumann
akan sama dengan peringkatnya sendiri,
\olref[ordinals][ordtype]{thmOrdinalRepresentation} kemudian akan memberi
tahu kita bahwa, setiap kali kita dapat mendeskripsikan suatu pengurutan baik
dalam teori himpunan kita, proses iteratif pembangunan himpunan harus
melampaui pengurutan baik tersebut.

Gagasan ini tentu tampak seperti konsekuensi dari \stagesinex. Sayangnya,
jika tujuan kita ialah memperoleh Penggantian dari gagasan ini, kita
menghadapi penghalang teknis yang sederhana: Penggantian benar-benar lebih
kuat daripada \olref[ordinals][ordtype]{thmOrdinalRepresentation}.
(Pengamatan ini dibuat oleh \citet[\S13.2]{Potter2004}; kita akan
membuktikannya dalam
\olref[sth][replacement][finiteaxiomatizability]{sec}.)

Akibatnya, jika kita hendak memahami \stagesinex{} dengan cara yang
menghasilkan Penggantian, prinsip itu tidak dapat \emph{sekadar} mengatakan
bahwa hierarki melampaui setiap pengurutan baik. Prinsip itu harus membuat
klaim yang lebih kuat. Untuk tujuan ini, \cite{Shoenfield:AST} mengusulkan
penguatan gagasan yang sangat alami sebagai berikut: hierarki tidak
\emph{kofinal} dengan himpunan mana pun.\footnote{G\"odel tampaknya telah
mengusulkan gagasan serupa; lihat \citet[p.~223]{Potter2004}. Untuk
pembahasan mengenai G\"odel dan \citeauthor{Shoenfield:AST}, lihat
\citet[90--5]{Incurvati2020}.} Secara sedikit lebih terperinci: jika $\tau$
merupakan pemetaan yang mengirim himpunan ke tahap-tahap hierarki, citra
sebarang himpunan $A$ di bawah $\tau$ tidak menghabiskan hierarki. Dengan
kata lain (secara skematis):
\begin{enumerate}
	\item[] \stagescofin. Jika $A$ merupakan himpunan dan $\tau(x)$ merupakan
	tahap bagi setiap $x \in A$, maka terdapat suatu tahap yang berada sesudah
	setiap $\tau(x)$ untuk $x \in A$.
\end{enumerate}
Jelas bahwa $\ZF$ membuktikan suatu versi \stagescofin{} yang diformalkan
secara tepat. Sebaliknya, kita dapat berargumen secara informal bahwa
\stagescofin{} membenarkan Penggantian.\footnote{Akan lebih sulit
membuktikan Penggantian dengan menggunakan suatu formalisasi
\stagescofin, karena $\Z$ sendiri tidak cukup kuat untuk mendefinisikan
tahap-tahap, sehingga tidak jelas bagaimana \stagescofin{} akan
diformalkan. Meskipun demikian, salah satu pilihannya ialah bekerja dalam
suatu perluasan $\LT$, sebagaimana dibahas dalam
\olref[sth][replacement][extrinsic]{sec}.} Andaikan
$(\forall x \in A)\lexists![y][\phi(x,y)]$. Maka, untuk setiap $x \in A$,
misalkan $\sigma(x)$ merupakan $y$ sedemikian sehingga $\phi(x,y)$, dan
misalkan $\tau(x)$ merupakan tahap ketika $\sigma(x)$ pertama kali terbentuk.
Berdasarkan \stagescofin, terdapat suatu tahap $V$ sedemikian sehingga
$(\forall x \in A)\tau(x)\in V$. Sekarang, karena setiap
$\tau(x) \in V$ dan $\sigma(x) \subseteq \tau(x)$, berdasarkan Pemisahan
kita dapat memperoleh
$\Setabs{y \in V}{(\exists x \in A)\sigma(x) = y} = \Setabs{y}{(\exists x \in A)\phi(x,y)}$.

\begin{prob}
	Formalkan \stagescofin{} di dalam $\ZF$.
\end{prob}

Jadi, \stagescofin{} membenarkan Penggantian. Dan sekurang-kurangnya masuk
akal bahwa \stagesinex{} membenarkan \stagescofin. Andaikan
\stagescofin{} gagal. Maka hierarki itu kofinal dengan suatu himpunan~$A$,
yakni kita mempunyai sebuah pemetaan $\tau$ sedemikian sehingga, untuk setiap
tahap~$S$, terdapat suatu $x \in A$ sehingga $S \in \tau(x)$. Dalam kasus
itu, kita mempunyai cara untuk memahami ``ketakhinggaan mutlak'' hierarki
yang diandaikan: hierarki itu \emph{dihabiskan} oleh jangkauan $\tau$ yang
diterapkan pada $A$. Dan hal itu melemahkan gagasan bahwa hierarki tersebut
``tak hingga secara mutlak''. Dengan mengambil kontraposisinya:
\stagesinex{} mengimplikasikan \stagescofin, yang selanjutnya membenarkan
Penggantian.

Ini merupakan upaya yang sungguh menjanjikan untuk memberikan pembenaran
intrinsik bagi Penggantian. Namun, apakah upaya itu akhirnya berhasil atau
tidak, harus kita serahkan kepada Anda untuk memutuskannya.

\end{document}
